\chapter{A centered angular packet closes the physical small-gap C2 chart}
\label{component:p7c5e0b890910}
\noindent\textbf{Component }\path{p7c5e0b890910}\par
{\small\noindent Source: \path{R13/advances/revision13_weak_local/centered_circuit_closure.tex}\par}
{\small\noindent SHA256: \path{7c5e0b8909104814fe60809f554c8de3ac83eb9d7c92370f9ac3b02a69b396ce}\par}
{\small\noindent Inventory use: WC1-WC28\par}
\medskip
\section{A centered angular packet closes the physical small-gap C2 chart}
\label{p7c5e0b890910:sec:r13-small-gap-circuit}

\begin{theorem}
Retain exactly the physical source hypotheses (BS1)--(BS3), the
original stationary contact, every strict source gap, and the original
positive infinitesimal total-scale intersection. Then the entire
original C2 parent is finite tensor SOS for every independent testing
pair with $u,v$ infinitesimal or zero, every sign, and every real or
complex state. In particular this includes $n=0$, both possibilities
for the bounded ratio $B/(1-r)$, and arbitrary relative rates of
$H/B$ and $1-r$.

Joining this local assertion to (GR2)--(GR7) and its printed exact
homogeneous endpoints gives an all-tests admission of precisely the
source chart (BS1). No other weak-C2 source-gap opening is included.
\end{theorem}

Write $a=z/B$, $t=y/B$, $h=H/B$, $e=1-r$. Thus $a,t$ are
positive bounded units with residues $\rho>0$, $\tau>1$, and
$B,e,h$ are positive infinitesimals. The event coefficients $n,m,C,A,D,
\alpha$ have their unchanged original meanings; $\lambda=\beta=0$.
All identities below are for $F/4$.

\paragraph{The angular packet that stays centered at every opening.}
The decisive source quantities are
\[
 q_s=nr^2+(m+\tfrac12)e^2>0,\qquad
 \widetilde m=m-a^2q_s.                                      \tag{WC1}
\]
The exact original source formulas, and not a comparison of rates,
give
\[
 \widetilde m/B\longrightarrow a_0^2t_0/2>0.                 \tag{WC2}
\]
More strongly $\widetilde m/B$ is jointly regular at
$B=h=0,r=1$: its denominator has a positive-unit residue there.
This remains true when $e/B$ is unbounded. The two pieces of $q_s$
must be kept together. Subtracting only the $nr^2$ piece leaves an
uncontrolled $e^2/B$ coefficient.

Here is a finite way to check all regularity assertions used below.
Put
\[
 c_a=1+aBr,\quad S_c=1+B(1-h+ar),\quad
 g=1+a(hB+r-1),\quad j=c_a-(1-h)/t,\quad \Delta=\delta/B.
\]
Use the original normalized formulas
\[
\begin{aligned}
 n&=\frac{(1+ar)c_aS_c-r\Delta}{2r\Delta},&
 m&=\frac{a^2(1+tB)e c_aS_c}{2g\Delta},\\
 C&=\frac{c_aM}{2(1-h)\Delta},&
 Ah&=\frac{t c_a(1+ar)S_c}{2\Delta},\\
 D&=\frac{(1+aB)(1+ar)(1-atBr)S_c}{2tj\Delta},&
 \alpha&=\frac{B^2(1+ar)M}{2(1-hB)\Delta},
\end{aligned}                                                   \tag{WC3}
\]
where
\[
 M=(1+aB)(1+ar)(h+t)-a(1+tB)e(1-h)
\]
and $\Delta$ is exactly the polynomial in (BS1).
Every denominator divides
\[
 \mathcal D=2r g(1-h)(1-hB)(t-1+h+aBtr)\Delta,
\]
whose residue is $2(t_0-1)(1+a_0)>0$. For an expression
linear in $n,m,C,Ah,D,\alpha$, clear $\mathcal D$ and its
printed pure-source denominator. Monic division of the numerator by
$B$, once or twice as printed, is exact. Evaluating the resulting
polynomial quotient at $B=h=0,r=1$ gives the residues below.
This supplies finite polynomial identities, rather than an appeal
to asymptotic contact cancellation. For (WC2), the constant numerator
vanishes identically at $B=0$ for every $h,r$; at that face
\[
 n=\frac e{2r(1+a)},\quad
 m=\frac{a^2e}{2(1-ae)(1+a)},\quad
 m/a^2=nr^2+(m+\tfrac12)e^2.
\]
The focused checker implements these finite divisions with
$\mathcal D$ kept as a common unit denominator.

\paragraph{Bounded testing ratio and all three hard state directions.}
Take $u,v>0$ and first suppose the pure-testing ratio $w=v/u$ is
bounded, including infinitesimal $w$. At $Z\ne0$ divide all states
by $Z$ and set $Z=1$; the homogeneous restoration is described below.
Set
\[
 E=tuW-v,\qquad N_y=(Y-huW)/\sqrt h,\qquad
 \xi=X+ru.                                                   \tag{WC4}
\]
Equivalently
\[
 W=(v+E)/(tu),\quad Y=h(v+E)/t+\sqrt h\,N_y,
 \quad X=\xi-ru.
\]
Only a pure-testing inverse and a pure-source square root occur.
In this frame the complete hard quadratic has bounded coefficients
and joint residue
\[
 (E,N_y,\xi):\qquad
 \operatorname{diag}\left(\frac1{2(t_0-1)},\frac{t_0}2,
                                      \frac12\right)\succ0.  \tag{WC5}
\]
For example $A=Ah/h$ always occurs multiplied by the square of
$h(v+E)/t+\sqrt h N_y$, so its normalized coefficients contain
$Ah$, $\sqrt h\,Ah$, or $hAh$, and no negative power of $h$.
The only testing-ratio dependence in the hard block is polynomial
in bounded $w$, such as $mw^2$; all its coefficients are bounded.
The displayed residue is obtained from the whole original four-state
parent, not a restricted soft graph.

\paragraph{Exact scalar identity with two retained weighted roots.}
Retain the scalar coefficients from (WA3)--(WA5), writing the
coefficient of $U W^2$ as $a_w$ to avoid collision with $a=z/B$:
\[
\begin{gathered}
 a_w=(C+D+\tfrac12)(1-h)^2+Ah\,h,\quad
 b_2=2mh+C(1-h)^2+Ah\,h+(\alpha+\tfrac12)h^2,\\
 b_4=h^2(m+\alpha+\tfrac12)+Ah\,h,\quad
 d_2=(n+\tfrac12)r^2+2Dr,\quad d_4=(n+D+\tfrac12)r^2,\\
 e_0=me+D(1-h)+\alpha h+\tfrac12,\quad
 e_1=mhe+r(D+\tfrac12)(1-h)+(\alpha+\tfrac12)h.
\end{gathered}                                                \tag{WC6}
\]
Let $x=aB$, $y=tB$, $k=\sqrt x$, $s=\sqrt{t/a}$ and
\[
 A_t=uv,\quad B_t=\frac{k v^2}{su},\quad C_t=xs,
 \quad G=kv,\quad F_v=V-y,\quad
 H_u=\frac a t\frac{v^2}u-u.
                                                               \tag{WC7}
\]
Then $A_tB_tC_t=G^3$, and all three circuit amplitudes equal $G$
at the actual contact $u=k,w=s$. Both $H_u$ and $F_v$ vanish
there. Define
\[
\begin{aligned}
 c_F&=(b_2-a^2nr^2)/t^2,\\
 q_V&=a_w/t^2+D-2e_0/t+2aq_s/t,\\
 f_A&=Ah\,h/t^2+d_2-2e_1/t+2anr^2/t+d_4U+b_4V/t^2,\\
 f_B&=\widetilde m/(ta^2B),\\
 f_C&=[\alpha-c_Fy^2+\alpha U]/(taB^2),\\
 f_G&=-(q_V+2c_Fy)/(aB).
\end{aligned}                                                    \tag{WC8}
\]
The entire scalar graph $E=N_y=\xi=0$ is exactly
\[
\begin{aligned}
 \mathcal S={}&q_sH_u^2+nr^2u^2H_u^2+c_FF_v^2\\
 &+f_AA_t^2+f_BB_t^2+f_CC_t^2-f_GG^2.             
\end{aligned} \tag{WC9}
\]
This is an unconditional polynomial identity after multiplying by
the pure-testing square $u^2$. To see its cancellation directly,
\[
 H_u^2=U-2aV/t+a^2V^2/(t^2U).
\]
Thus $q_sH_u^2$ removes the full $q_sU$ coefficient and replaces
$m$ by $\widetilde m$ in the $V^2/U$ coefficient. The second
square removes $nr^2U^2$ and changes only the displayed $V^2$ and
$UV$ coefficients. Nothing has been discarded from the original
parent, and source physicality is not transferred to a remainder.

The finite source divisions described after (WC3) give
\[
 \begin{gathered}
 c_F\to1/(2t_0),\quad f_A,f_B,f_C\to\tfrac12,
 \quad f_G\to\tfrac32,\\
 q_V/B\to-1-3a_0/2,\qquad
 \partial_U f_C\to(1+a_0)/(2a_0).                  
 \end{gathered} \tag{WC10}
\]
All coefficient arrays in $f_A,f_C$ are bounded; each is affine
in $U,V$. The constant $f_B,f_G$ have no testing dependence.
These are joint regular reductions and contain no $e/B$ or $h/e$
assumption. The two positive weighted roots
\[
 R_q=\sqrt{q_s}\,H_u,\qquad R_n=\sqrt n\,r uH_u            \tag{WC11}
\]
remain inside the root list, including literal $n=0$.

\paragraph{Finite contact lifting of the scalar circuit.}
Apply the already proved three-circuit contact-lift lemma with
$\lambda=1/2$ to the last four terms of (WC9), retaining the strict
$c_FF_v^2$ pivot and the two squares (WC11). Its conductors are
the exact separated identities
\[
 G(u-k)=k(A_t-G),\qquad G(w-s)=s(B_t-G).                \tag{WC12}
\]
The required finite Taylor arrays are precisely (IC11)--(IC14),
with the present $f_A,f_B,f_C,f_G$ and $x,y,s$. They apply because
these functions are affine in $U=u^2,V=u^2w^2$, with bounded
pure-source coefficients. Actual source contact gives all three
value/first-jet balances: the squares removed in (WC9) have zero
value and first derivative there. No Hessian of a modified physical
parent is asserted. Consequently the scalar has a finite strict
bounded-residue Gram on $R_q,R_n,F_v$ and the finite three-circuit
root list.

\paragraph{Every full first normal is allocated on the same roots.}
Use the complete graph normal coefficients $f_w,f_z,k_w,k_z$ in
(WA6)--(WA7), equivalently
\[
\begin{aligned}
 f_w&=-m+[-mh+(D+\tfrac12)(1-h)]U,\\
 f_z&=-nr(1+U)(1+V)+me-(D+\tfrac12)V(1+rU)
                                      +\tfrac12e(1+V),\\
 k_w&=-(C+D+\tfrac12)(1-h)+Ah+Ah\,U\\
 &\quad+[-C(1-h)+Ah+m+(\alpha+\tfrac12)h]V
                  +[Ah+h(m+\alpha+\tfrac12)]UV,\\
 k_z&=D-\alpha+[-me+r(D+\tfrac12)-\alpha-\tfrac12]U.
\end{aligned}                                                    \tag{WC13}
\]
Put
\[
 P_Y=2(k_w/t+k_z),\qquad
 P_E=2(A_w/U-mV/U)/t^2-2(e_0+e_1U)/t,
\]
where $A_w$ is the complete $W^2$ coefficient in (WS2).
The exact first derivatives of $F/4$ in the frame (WC4) are
\[
\begin{aligned}
 \mathcal L_{N_y}&=\sqrt h\,vP_Y,\\
 \mathcal L_E&=2mvw^2/t^2+vP_E,\\
 \mathcal L_\xi&=2u\{w^2f_w/t+f_z\}.
\end{aligned}                                                    \tag{WC14}
\]
These identities are verified by differentiating the entire original
parent after (WC4), retaining all original packets.

Extract the following complete angular contributions:
\[
 \begin{aligned}
 \mathcal L_\xi&=-2a\sqrt{q_s}\,R_q+2\sqrt n\,uR_n+\mathcal R_X,\\
 \mathcal L_E&=2a\sqrt{q_s}\,(w/t)R_q+\mathcal R_E,\\
 \mathcal L_{N_y}&=\sqrt h\,\mathcal R_Y,
 \quad \mathcal R_Y=vP_Y.
 \end{aligned}                                                   \tag{WC15}
\]
The displayed coefficients are infinitesimal because $q_s,n$ are
infinitesimal and $u,w$ are bounded. No inverse of either weight
appears. The exact residuals in (WC15) are finite polynomials in
$u,w$ of the following explicitly checkable form:
\[
 \mathcal R_\bullet
 =B\sum_{j=0}^{3}c_{\bullet j}u w^j
   +\sum_{i\ge3,\,j\ge1}d_{\bullet ij}u^iw^j,       \tag{WC16}
\]
where all source coefficients are bounded; only the monomials
actually occurring in (WC13)--(WC15) are included. Their exact
supports are
\[
\begin{array}{c|c|c}
 &Buw^j&u^iw^j\ (i\ge3)\\ \hline
 X&j=0,2&(i,j)=(3,2),(5,2)\\
 E&j=1,3&(i,j)=(3,1),(3,3),(5,3)\\
 Y&j=1&(i,j)=(3,1),(3,3),(5,3).
\end{array}                                                    \tag{WC17}
\]
Equation (WC16) is a finite numerator-division assertion, checked
with (WC3) and the common unit denominator $\mathcal D$. It does
not follow merely from normal vanishing. In particular the apparent
uncentered $e/B$ terms cancel in the $u w^0,u w^2,u w^3$ entries
only after retaining all of $q_s$ in (WC15).

For clarity, its complete arrays may be written without any
coefficient-extraction convention. Put
\[
\begin{aligned}
 K_X&=-nr+(m+\tfrac12)e-aq_s,\\
 K_E&=a_w/t^2-e_0/t+aq_s/t,\\
 K_Y&=[-(C+D+\tfrac12)(1-h)+Ah]/t+D-\alpha,\\
 x_3&=2\{(D+\tfrac12)[(1-h)/t-1]+e/2
                             -nr(1+a/t)-mh/t\},\\
 x_5&=-2r(n+D+\tfrac12).
\end{aligned}
\]
Then exactly
\[
\begin{aligned}
 \mathcal R_X={}&2K_Xu-2\widetilde m\,uw^2/t
                                  +x_3u^3w^2+x_5u^5w^2,\\
 \mathcal R_E={}&2K_Euw+2\widetilde m\,uw^3/t^2
 +(2Ah\,h/t^2-2e_1/t)u^3w
                    +2b_2u^3w^3/t^2+2b_4u^5w^3/t^2,\\
 \mathcal R_Y={}&2K_Yuw
 +2[Ah/t-me+r(D+\tfrac12)-\alpha-\tfrac12]u^3w\\
 &+2[-C(1-h)+Ah+m+(\alpha+\tfrac12)h]u^3w^3/t\\
 &+2[Ah+h(m+\alpha+\tfrac12)]u^5w^3/t.
\end{aligned}                                                    \tag{WC17a}
\]
The only source quotients required for these arrays are
$\widetilde m/B,K_X/B,K_E/B,K_Y/B$. Each has a unit denominator
after exact monic division. Their numerators vanish at $B=0$
for every $h,r$, not merely at the joint corner. For example
$m/a^2=q_s$ and $-nr+(m+1/2)e=m/a$ on that face prove
$K_X=0$. Substituting (WC3) at $B=0$ gives $K_Y=0$ directly;
that equality implies $a_w=Ah+t(1-h)D$, which reduces $K_E=0$
to $Ah/t-me-1/2+m/a=0$. These are exact face identities with
the same unit denominators, so the four polynomial divisions
prove (WC16) for arbitrary independent source openings.

Every monomial in (WC16) now has an explicit legal circuit allocation:
\[
\begin{array}{rclcrcl}
 Bu&=&uC_t/(as),&&Buw&=&(B/k)G,\\
 Buw^2&=&(Bs/k)B_t,&&Buw^3&=&(Bs/k)wB_t,\\
 u^iw^j&=&u^{i-2}w^{j-1}A_t&&(i\ge3,j\ge1).
\end{array}                                                    \tag{WC18}
\]
All resulting coefficients are infinitesimal. Their finite first
divided differences in $u$ are bounded, and the divided differences
in $w$ after setting $u=k$ are infinitesimal: every $w$-dependent
coefficient retains $u$ or $k$. Both properties follow directly
from the displayed monomials, with no testing-ratio inverse.
Each residual normal vanishes at the actual contact, because the
original complete normal and $H_u$ vanish there. The monic division
(IC21), followed by (WC12), therefore expresses each full normal
with infinitesimal coefficients on the same strict circuit roots.
Together with (WC5), (WC9), and (WC15), this proves a finite strict
bounded-residue Gram for the entire parent at every bounded $w$.

\paragraph{The reciprocal ratio uses the physical source scale bound.}
Suppose $w=v/u$ is unbounded, so $\omega=u/v$ is infinitesimal.
Return to homogeneous states. Set
\[
 p=uW,\quad q=vZ,\quad E=tp-q,\quad
 Y=hp+\sqrt hN_y,\quad \xi=X+ruZ,
\]
and use the complete root list
\[
 E,\ N_y,\ \xi,\quad P=\sqrt m\,vW,\quad
 Q=\sqrt\alpha Z,\quad R=u^2W,\quad S=uvW.          \tag{WC19}
\]
Here $m/e$ and $\alpha/B^2$ are positive bounded units. BS3 gives
$B/e$ and $n/e$ bounded, hence $n/m$ bounded. The pure-source
comparison coefficients
\[
 C_*=\frac{t+h}{2(1-h)},\qquad (Ah)_*=t/2,\qquad
 D_*=\frac1{2(t+h-1)}
\]
satisfy
\[
 C-C_*=B c_B+e c_e,\quad Ah-t/2=B a_B+e a_e,
 \quad D-D_*=B d_B+e d_e,                            \tag{WC20}
\]
with bounded pure-source coefficients. To check this finitely,
clear the unit denominators in (WC3), subtract the numerator at
$B=0$, divide that difference by $B$, and divide its remaining
$B=0$ numerator by $r-1=-e$. The latter numerator vanishes
identically at $r=1$ for every $h$. These are the exact finite
divisions implemented by the focused checker.

The quadratic comparison base
\[
 C_*(p-Y)^2+\frac{t}{2h}Y^2+D_*(p-Y-q)^2
                         +\tfrac12(p-Y)^2-pq
\]
is exactly
\[
 D_*E^2-2D_*\sqrt h\,EN_y+
 \frac{t(t+2h-1)}{2(1-h)(t+h-1)}N_y^2.              \tag{WC21}
\]
It has strict diagonal residue $1/[2(t_0-1)],t_0/2$.
The comparison error has infinitesimal coefficients on (WC19):
\[
 \sqrt e\,p=\sqrt{e/m}\,\omega P,\qquad q=tp-E,
 \qquad Y=hp+\sqrt hN_y.                            \tag{WC22}
\]
For the $A$ term use the exact expansion
$(Ah-t/2)Y^2/h=(Ah-t/2)
[hp^2+2\sqrt h\,pN_y+N_y^2]$. Thus no $e/h$ or $h/e$
comparison is introduced.

With $X+uZ=0$, the exact higher-angle part of the $C,A,D$ packets
and half core, after removing their quadratic base, is
\[
\begin{aligned}
 &U\{AY^2-2Dq(p-Y-q)+q^2/2-pq\}\\
 &\quad+V\{C(p-Y)^2+(A+\tfrac12)Y^2\}
 +(A+\tfrac12)UVY^2+(D+\tfrac12)U^2q^2.             
\end{aligned} \tag{WC23}
\]
Substitute (WC22). Its residue on $(R,S)$ is strictly diagonal,
with pivots $t_0^2/2,t_0/2$. Every other coefficient is bounded
and infinitesimal: for instance $UpN_y=uRN_y$, $VpE=vSE$,
$UVp^2=US^2$. Terms with $A=Ah/h$ are regular by the same
explicit $Y$ expansion used after (WC22).

To restore the complete $X$ dependence write
\[
 B_x=X+uZ=\xi+euZ
 =\xi+t(e/\sqrt m)\omega^2P-e\omega E.              \tag{WC24}
\]
Both added coefficients are infinitesimal. Relative to $B_x=0$,
the exact $D$ packet and half core add
\[
 2(D+\tfrac12)uvB_x(p-Y-(1+U)q)
 +(D+\tfrac12)UVB_x^2+\tfrac12(1+V)B_x^2.           \tag{WC25}
\]
The last term supplies the strict $\xi$ pivot $1/2$; all remaining
Gram coefficients are infinitesimal, using $uvp=uS$.
The remaining three original positive packets retain their full roots:
\[
\begin{aligned}
 \sqrt m\,M&=P-\sqrt m\,B_x+\sqrt m\,uvY,\\
 \sqrt\alpha\,(vY-Z)&=-Q+\sqrt\alpha hS
                                      +\sqrt{\alpha h}\,vN_y,\\
 \sqrt n\,X&=\sqrt n\,\xi+r\sqrt n\,\omega E
                                  -rt\sqrt{n/m}\,\omega^2P.
\end{aligned}                                                    \tag{WC26}
\]
Retain the original positive pure-testing $D_u,D_v$ factors.
The first two roots have leading entries $P,-Q$; every coefficient
of the third is infinitesimal, including $n=0$. No square or
first normal was dropped. Hence the complete parent has a bounded
Gram on (WC19) with strict diagonal residue
\[
 \operatorname{diag}\left(\frac1{2(t_0-1)},\frac{t_0}2,
      \frac12,1,1,\frac{t_0^2}2,\frac{t_0}2\right)\succ0.
                                                               \tag{WC27}
\]
This proves the reciprocal ratio without comparing any source scale
with $u$ or $v$.

\paragraph{Restoration, verification, and exact admission.}
The exact sign gauges $(u,Y,Z)\mapsto(-u,-Y,-Z)$ and
$(v,X,Z)\mapsto(-v,-X,-Z)$ preserve the entire original parent,
including the $nX^2$ packet. Literal axes have the exact identities
\[
\begin{aligned}
 (F/4)|_{u=0}={}&nD_vX^2+m(vW-X)^2
 +(C+A+\tfrac12)D_vY^2+D(Y+vZ)^2\\
 &+\alpha(vY-Z)^2+\tfrac12D_vX^2,\\
 (F/4)|_{v=0}={}&nD_uX^2+(m+\tfrac12)(X+uZ)^2
 +(C+D+\tfrac12)(uW-Y)^2\\
 &+AD_uY^2+\alpha D_uZ^2.
\end{aligned}                                                    \tag{WC28}
\]
All displayed coefficients are nonnegative. Homogenize the
bounded-ratio proof by multiplying its
identity by $Z^2$: $E,N_y,\xi$ become the homogeneous linear
forms, and each scalar circuit root and $H_u$ is multiplied by $Z$.
No $Z$ inverse remains, so the identity extends to $Z=0$.
The same real identity applied to the real and imaginary state parts
gives the complex-state assertion. The previously proved GR join
supplies every remaining independent testing chart and projective
endpoint on this exact source chart.

The new checker verifies the scalar identity (WC9), all three
complete first normals by direct original-parent differentiation,
every entry of the hard residue, the finite source divisions and
normal supports, the exact reciprocal higher-angle identity, and
both strict higher-angle pivots. These are focused checks of the
new argument; no previous verification suite is replayed. The
finite Gram/positivity proof is the argument above. It closes only
the BS1 small-$1-r$, $z/B$-unit source chart, including its printed
strict gaps and $n=0$; the other C2 source-gap complements remain
separate obligations.
